\documentclass[10pt,letterpaper]{article}
\usepackage[margin=0.6in]{geometry}
\usepackage[T1]{fontenc}
\usepackage{lmodern}
\usepackage{enumitem}
\usepackage{titlesec}
\usepackage[hidelinks]{hyperref}
\setlength{\parindent}{0pt}
\setlength{\parskip}{3pt}
\setlist[itemize]{nosep,leftmargin=1.3em}
\titleformat{\section}{\large\bfseries}{}{0pt}{}
\titlespacing*{\section}{0pt}{12pt}{5pt}
\newcommand{\entry}[2]{\par\textbf{#1}\hfill #2\par}
\newcommand{\paper}[3]{\par\textbf{#1}\par #2\par #3\par}
\begin{document}
{\LARGE\bfseries Hansol Suh}\hfill Last updated: September 14, 2026\par
\href{https://hansolsuh.com}{hansolsuh.com} \quad
\href{mailto:hansol.suh.123@gmail.com}{hansol.suh.123@gmail.com}\par
\href{https://github.com/hansolsuh}{GitHub: hansolsuh} \quad
\href{https://gitlab.com/hsuh}{GitLab: hsuh}

\section{Education}
\entry{Georgia Institute of Technology}{Atlanta, GA}
\entry{MSc, Computational Science and Engineering}{Fall 2022}
PhD studies, Computational Science and Engineering: Spring 2018--Fall 2022.
Program concluded with an MSc. Advisor: Tobin Isaac. Research: numerical optimization and HPC.
\entry{BSc, Civil Engineering, Highest Honor}{August 2011--July 2016}

\section{Technical Skills}
\textbf{Programming:} C, C++, Python, Julia, Java, MATLAB; CUDA and HIP.\par
\textbf{Methods and tools:} MPI, OpenMP, Unix/Linux, Git, machine learning, parallel computing.

\section{Publications}
\textbf{Peer-reviewed publications}
\paper{Verification of PETSc with CIVL Using LLM-Generated ACSL Contracts and Deterministic Driver Generation}
{\textbf{H. Suh}, J. H\"uckelheim, and S. Siegel.}
{Correctness 2026: 10th International Workshop on Software Correctness for HPC Applications (SC26). Accepted; workshop scheduled for November 16, 2026.\par
\href{https://arxiv.org/abs/2609.31687}{arXiv:2609.31687}}
\paper{PETSc/TAO Developments for GPU-Based Early Exascale Systems}
{R. T. Mills, M. F. Adams, S. Balay, J. Brown, J. Faibussowitsch, T. Isaac, M. G. Knepley, T. Munson, \textbf{H. Suh}, S. Zampini, H. Zhang, and J. Zhang.}
{The International Journal of High Performance Computing Applications. 2025.\par
\href{https://doi.org/10.1177/10943420241303710}{doi:10.1177/10943420241303710}}
\paper{Evaluation of a Minimally Synchronous Algorithm for 2:1 Octree Balance}
{\textbf{H. Suh} and T. Isaac.}
{SC20: International Conference for High Performance Computing, Networking, Storage and Analysis. November 2020.}

\textbf{Posters}
\paper{Discrete Adjoints in PFLOTRAN for Physics-Constrained Machine Learning}
{\textbf{H. Suh} et al.}
{AGU Annual Meeting 2026, San Francisco, CA. Submitted.}

\textbf{Preprints and manuscripts}
\paper{A Portable and Versatile Limited-Memory BFGS Implementation in PETSc/TAO}
{\textbf{H. Suh}, T. Isaac, A. Dener, T. Munson, H. Zhang, and R. T. Mills.}
{\href{https://arxiv.org/abs/2608.00196}{arXiv:2608.00196}. 2026. Submitted to \emph{Computational Optimization and Applications}.}
\paper{Variable Metric Three Operator Splitting}{\textbf{H. Suh} and T. Isaac.}
{\href{https://arxiv.org/abs/2610.38935}{arXiv:2610.38935}. 2026.}
\paper{A Heuristic for Parameter Selection in Adaptive Relaxed ADMM Based on Numerical Ranges}
{T. Isaac and \textbf{H. Suh}.}{Manuscript.}

\newpage
\section{Technical Reports and Manuals}
\paper{PETSc/TAO Users Manual, Revision 3.25}
{S. Balay et al. (including \textbf{H. Suh}).}
{ANL-21/39 Rev. 3.25. Argonne National Laboratory. 2026.}
\paper{TAO Users Manual, Revision 3.15}
{A. Dener, A. Denchfield, \textbf{H. Suh}, T. Munson, J. Sarich, S. Wild, S. Benson, and L. C. McInnes.}
{ANL/MCS-TM-322 Rev. 3.15. Argonne National Laboratory. March 2021.}

\section{Experience}
\entry{Argonne National Laboratory}{Lemont, IL}
\entry{Predoctoral Appointee}{March 2023--Present}
\begin{itemize}
\item Develop dense BFGS implementations for GPUs in PETSc/TAO.
\item Develop proximal optimization solvers for PETSc/TAO.
\item Supervisors: Tobin Isaac (former); Richard Mills (current).
\end{itemize}
\entry{Summer Research Appointments}{2019, 2021, 2022}
\begin{itemize}
\item Research Aide, Technical (June--August 2022): developed Anderson-accelerated three-operator splitting and hierarchical, composable PETSc/TAO solvers. Supervisor: Todd Munson.
\item Research Aide (May--August 2021): implemented variable-metric three-operator splitting in PETSc/TAO and evaluated the implementation in an HPC setting. Supervisor: Alp Dener.
\item Givens Associate (May--August 2019): implemented ADMM and a variant in PETSc/TAO and evaluated their performance in an HPC setting. Supervisor: Alp Dener.
\end{itemize}
\entry{Georgia Institute of Technology}{Atlanta, GA}
\entry{Graduate Research Assistant}{May 2018--December 2022}
Research with Tobin Isaac in PETSc development, numerical optimization and analysis, p4est, and HPC.
\entry{Graduate Student}{August 2016--May 2018}
\entry{Republic of Korea Army}{Icheon, South Korea}
Anti-Aircraft Gunner; Armory Department.\hfill November 2013--August 2015

\section{Grants \& Funding}
\entry{CS2: Modular Verification of Scientific Software}{May 2025--May 2029}
\textbf{Using Model Checking, Symbolic Execution, and Abstraction}\par
Senior Personnel, NSF/DOE Correctness for Scientific Computing Systems (CS2), NSF 24-571.
Award: \$3,200,000. PIs: Stephen Siegel (University of Delaware), Jan H\"uckelheim (Argonne National Laboratory).
\entry{AI4Soil}{Summer 2026}
A Deployable Artificial Intelligence Soil Prototype for Improving Land--Atmosphere Prediction.
ANL LDRD Seed project; collaborator with Joshua Minai. Award: \$25,000.

\section{Projects}
\entry{Terrestrial Wetland Function and Resilience Science Focus Area}{Summer 2025--Present}
Developing adjoint-based sensitivity-analysis capabilities in PFLOTRAN for subsurface reactive transport, and physics-constrained machine-learning methods for wetland systems.
\par\textbf{TaoTerm}\par
PIONEER thrust: Computationally Efficient Composable Solvers for PETSc/TAO.
Developed composable objects for the TAO optimizer, enabling flexible, on-the-fly modification of optimization objectives.

\newpage
\section{Presentations}
\paper{TaoTerm: Composable Objective Function Terms}{\textbf{H. Suh}.}
{Consortium for the Advancement of Scientific Software BoF Days, online. February 11, 2026.\par
\href{https://hansolsuh.com/assets/taoterm_2026_bof_cass-4.pdf}{Slides}}
\paper{Dense Broyden--Fletcher--Goldfarb--Shanno (BFGS)}{\textbf{H. Suh}.}
{PETSc User Meeting 2025, Buffalo, NY. May 20, 2025.\par
\href{https://hansolsuh.com/assets/petsc2025_bfgs_slides.pdf}{Slides}}
\paper{An Introduction to PETSc/TAO}{\textbf{H. Suh}.}
{ALCF Hands-On HPC Workshop, Lemont, IL. October 31, 2024.\par
\href{https://hansolsuh.com/assets/alcf_workshop_2024_petsc.pdf}{Slides}}
\paper{Mini Tutorial: TaoADMM}{\textbf{H. Suh}.}
{PETSc User Meeting 2023, Chicago, IL. June 7, 2023.\par
\href{https://hansolsuh.com/assets/taoadmm.pdf}{Slides}}
\paper{An Asynchronous Algorithm for 2:1 Octree Balance}{\textbf{H. Suh} and T. Isaac.}
{SIAM Conference on Parallel Processing 2020, Seattle, WA. February 14, 2020.}
\paper{Using the PETSc/TAO ADMM Methods on GPUs}{\textbf{H. Suh}, A. Dener, T. Isaac, and T. Munson.}
{SIAM Conference on Parallel Processing 2020, Seattle, WA. February 13, 2020.}

\section{Teaching}
\entry{Computational Data Analysis (ISYE 6740)}{Fall 2021, Fall 2022}
Teaching Assistant.
\entry{Introduction to High Performance Parallel Computing (CX 4220)}{Spring 2021}
Teaching Assistant.
\entry{High Performance Parallel Computing (CSE 6230)}{Fall 2019}
Teaching Assistant.
\end{document}
